\documentclass[10pt]{article}

\usepackage[utf8]{inputenc}

\usepackage[T1]{fontenc}

\usepackage{amsmath}

\usepackage{amsfonts}

\usepackage{amssymb}

\usepackage[version=4]{mhchem}

\usepackage{stmaryrd}

\usepackage{graphicx}

\usepackage[export]{adjustbox}

\graphicspath{ {./images/} }

\usepackage{bbold}



\begin{document}

д е л е н и е м) С. с. $K$. Клеточные пространства $|\mathrm{Bd} K|$ и $|K|$ естественно гомеоморфны (но не изоморфны). При этом гомеоморфизме каждая верпина $|s|$ из $|\mathrm{Bd} K|$ (т. ө. нульмерная клетка, отвечающая вершине $s \in \mathrm{Bd} K$ ) переходит в центр тяжести (барицентр) замкнутого симплекса $|s| \subset|K|$.



C. с. $\mathrm{Bd} K$ естественным образом упорядочена. Если C. с. $K$ упорядочена, то соответствие $s \rightarrow$ (первая вершина $s$ ) определяет симплициальное отображение $\mathrm{Bd} K \rightarrow K$, сохраняющее упорядоченность. Оно наз. к а но ни ч е с ки м с д в и го м. Его геометрич. реализация (являющаяся непрерывным отображением $|\mathrm{Bd} K| \rightarrow|K|)$ гомотопна естественному гомеоморфизму $|\mathrm{Bd} K| \rightarrow|K|$.



Симплициальное отображение $\varphi: K \rightarrow L$ (или его геометрич. реализация $|\varphi|:|K| \rightarrow|L|)$ наз. с и м п л иц и а ль ь о й а п п р о к с и м а ци е й непрерывного отображения $f:|K| \rightarrow|L|$, если для каждой точки $\alpha \in|K|$ точка $|\varphi|(\alpha)$ принадлежит минимальному замкнутому симплексу, содержащему точку $f(\alpha)$, т. е., что равносильно, если для каждой вершины $x \in K$ имеет место вложение $f(\operatorname{St} x) \subset \operatorname{St} \varphi(x)$. При әтом отображения $f$ и $|\varphi|$ гомотопны.



Те о рема осимплици а льной а ппрокс и м а ц и и утверждает, что если С. с. $K$ конечна, то для каждого непрерывного отображения $f:|K| \rightarrow|L|$ найдется такое число $N$, что для всех $n \geqslant N$ существует симплициальная аппроксимация $\mathrm{Bd}^{n} K \rightarrow L$ отображения $f$ (рассматриваемого как отображение $\left|\mathrm{Bd}^{n} K\right| \rightarrow$ $\rightarrow|L|)$.



Лит.: [1] С п е н ь е р Э., Алгебраическая топология, пер. с англ., М., 1971; [2] Х и Л т о н П. Д ж., У а й Л и С., Теория гомологий. Введение в алгебраическую топологию, пер. с англ., Soc.», 1939, v. 45 , p. 243-327. СИМПЛИЦАЛЬНОЕ МНОЖЕСТВО (прежние названия - по л у с и м п лици а льны й ко мпле к с, полны й полусимплици а льны й к о м п Л е к с) - симплициальный объект категории множеств Ens, т. е. система множеств ( $n-\mathrm{x}$ с ло е в) $K_{n}, n \geqslant 0$, связанных отображениями $d_{i}: K_{n} \rightarrow K_{n-1}$, $0 \leqslant i \leqslant n \quad\left(\right.$ о п р р а т о р а м и г р а н е й), и $s_{i}: K_{n} \rightarrow$ $\rightarrow K_{n+1}, \quad 0 \leqslant i \leqslant n \quad($ п е р а т о р а м и в ы р о ж д ен и я), удовлетворяющих соотношөниям



\[

\begin{aligned}

& d_{i} d_{j}=d_{j-1} d_{i}, \text { если } i<j,

\end{aligned}

\]



\begin{center}

\includegraphics[max width=\textwidth]{2023_07_08_cc11d7597b9578fc56d9g-1}

\end{center}



Точки слоя $K_{n}$ наз. $n$-м е р ны м и с и м п ле к с а м и C. м. $K$. Если заданы только операторы $d_{i}$, удовлетворяющие соотнопшениям $d_{i} d_{j}=d_{j-1} d_{i}, i<j$, то система $\left\{K_{n}, d_{n}\right\}$ наз. по л у с и м п ли и и а льны м множ е с т в о м.



Симплици альным отображением $f: K \rightarrow K^{\prime}$ С. м. $K$ в С. м. $K^{\prime}$ наз. морфизм функторов, т. е. последовательность отображений $f_{n}: K_{n} \rightarrow K_{n}^{\prime}$, $n \geqslant 0$, удовлетворяющих соотношениям



\[

d_{i} f_{n+1}=f_{n} d_{i}, \quad 0 \leqslant i \leqslant n+1 ; \quad s_{i} f_{n}=f_{n+1} s_{i}, 0 \leqslant i \leqslant n .

\]



С. м. и их симплициальные отображения образуют категорию $\Delta^{0}$ Ens. Если все отображения $f_{n}$ являются вложениями, то С. м. $K$ наз. с и м п ли ц и а ль ны м п о д м н о же с т в о м С. м. $K^{\prime}$. При этом операторы граней и вырождения в С. м. $K$ представляют собой ограничения соответствующих операторов в C. M. $K^{\prime}$. Для любого топологич. пространства $X$ определено С. м. $S(X)$, наз. с и н г у л я р н ы м С. м. пространства $X$, симплексами к-рого являются сингулярные симплексы пространства $X$ (см. Сингулярные гомологии), т. е. непрерывные отображения $\sigma: \Delta^{n} \rightarrow X$, где $\Delta^{n}-$ $n$-мерный геометрический стандартный силплекс



\[

\Delta^{n}=\left\{\left(t_{0}, \ldots, t_{n}\right) \mid 0 \leqslant t_{i} \leqslant 1, \sum_{i=0}^{n} t_{i}=1\right\} \subset \mathbb{R}^{n+1} .

\]



Операторы граней $d_{i}$ и вырождения $s_{i}$ әтого С. м. определяются формулами



$\left(d_{i} \sigma\right)\left(t_{0}, \ldots, t_{n-1}\right)=\sigma\left(t_{0}, \ldots, t_{i-1}, 0, t_{i}, \ldots, t_{n-1}\right)$, $\left(s_{i} \sigma\right)\left(t_{0}, \ldots, t_{n+1}\right)=$



\[

=\sigma\left(t_{0}, \ldots, t_{i-1}, t_{i}+t_{i+1}, t_{i+2}, \ldots, t_{n+1}\right) .

\]



Соответствие $X \mapsto S(X)$ является функтором (наз. с и н г у л я р ны м ф ун к то ро м) из категории топологич. пространств Тор в категорию С. м. $\Delta^{0}{ }^{0} \mathrm{Ens}$.



Произвольная симплициальная схема $K$ определяет С. м. $O(K)$, симплексами размерности $n$ к-рого являются $(n+1)$-членные последовательности $\left(x_{0}, \ldots, x_{n}\right)$ вершин схемы $K$, обладающие тем свойством, что в $K$ существует такой симплекс $s$, что $x_{i} \in s$ для любого $i=0$, $1, \ldots, n$. Операторы граней $d_{i}$ п вырождения $s_{i}$ этого С. м. определяются формулами



\[

d_{i}\left(x_{0}, \ldots, x_{n}\right)=\left(x_{0}, \ldots, \hat{x}_{i}, \ldots, x_{n}\right),

\][



s\_\{i\}\textbackslash left(x\_\{0\}, ···, x\_\{n\}\textbackslash right)=\textbackslash left(x\_\{0\}, ···, x\_\{i\}, x\_\{i\}, x\_\{i+1\}, ···, x\_\{n\}\textbackslash right),

]



где знак $\wedge$ означает, что символ, стоящий под ним, опускается. Если симплициальная схема $K$ упорядочена, то симплексы $\left(x_{0}, \ldots, x_{n}\right)$, для к-рых $x_{0} \leqslant \ldots \leqslant x_{n}$, образуют симплициальное подмножество $O^{+}(K)$ С. м. $O(K)$. Соответствие $K \mapsto O(K)$ (и $\left.K \longmapsto O^{+}(K)\right)$ является функтором из категории симплициальных схем (упорядоченных симплициальных схем) в категорию $\Delta^{0}$ Ens. Для произвольной группы $\pi$ определено С. м. $K(\pi)$, симплексами размерности $n$ к-рого являются классы пропорциональных $(n+1)$-членных последовательностей $\left(x_{0}, \ldots, x_{n}\right), x_{i} \in \pi$ (по определению $\left(x_{0}, \ldots, x_{n}\right)$ $\sim\left(x_{0}^{\prime}, \ldots, x_{n}^{\prime}\right)$, если существует такой элемент $y \in \pi$, что $x_{i}^{\prime}=y x_{i}$ для всех $\left.i=0, \ldots, n\right)$. Операторы граней $d_{i}$ и вырождсния $s_{i}$ С. м. $K(\pi)$ определяются формулами



\[

\begin{gathered}

d_{i}\left(x_{0}: \ldots: x_{n}\right)=\left(x_{0}: \ldots: x_{i-1}: x_{i+1}: \ldots: x_{n}\right), \\

s_{i}\left(x_{0}: \ldots: x_{n}\right)=\left(x_{0}: \ldots: x_{i-1}: x_{i}: x_{i}: x_{i+1}: \ldots: x_{n}\right) .

\end{gathered}

\]



C. м. $K(\pi)$ является на самом деле симплициальной группой.



Для произвольной абелевой групшы $\pi$ и любого целого $n \geqslant 1$ определено С. м. (на самом деле, симплициальная абелова группа) $E(\pi, n)$, симплексами размерности $q$ к-рого являются $n$-мерные коцепи $q$-мерного геометрического стандартного симплекса $\Delta^{q}$ с коэффициентами в группе $\pi$ (таким образом, $E(\pi, n)_{q}=$ $\left.=C^{n}(\Delta q ; \pi)\right)$. Обозначая вершины симплекса $\Delta^{q}$ символами $e_{j}^{q}, j=0, \ldots \dot{q}$, определяют симплициальные отображения $\delta_{i}: \Delta^{q-1} \rightarrow \Delta^{q}$ и $\sigma_{i}: \Delta^{q} \rightarrow \Delta^{q-1}$ формулами



\[

\begin{aligned}

& \delta_{i}\left(e_{j}^{q-1}\right)=\left\{\begin{array}{l}

e_{j}^{q}, \quad \text { если } j<i, \\

e_{j+1}^{q}, \text { если } j \geqslant i ;

\end{array}\right. \\

& \sigma_{i}\left(e_{j}^{q}\right)=\left\{\begin{array}{c}

e_{j}^{q-1}, \text { если } j \leqslant i, \\

e_{j-1}^{q-1}, \text { если } j>i .

\end{array}\right.

\end{aligned}

\]



Индуцированные гомоморфизмы групп коцепей



\[

\begin{aligned}

d_{i} & =\delta_{i}^{*}: C^{n}\left(\Delta^{q} ; \pi\right) \longrightarrow C^{n}\left(\Delta^{q-1} ; \pi\right), \\

s_{i} & =\sigma_{i}^{*}: C^{n}\left(\Delta^{q-1} ; \pi\right) \longrightarrow C^{n}\left(\Delta^{q} ; \pi\right)

\end{aligned}

\]



являются, по определению, операторами граней п





\end{document}